\section{Sanctity, Legacy, and Unresolved Questions}

Pope Saint Gregory the Great stands at the exact hinge between Christian
antiquity and the Latin Middle Ages. He was the last great Latin Father formed
inside the classical senatorial aristocracy of ancient Rome, and the first
bishop of Rome who had to imagine a Christian West that no longer depended on
an emperor at Constantinople. Unlike Augustine, he did not forge vast
speculative architectures of philosophy and grace; unlike Jerome, he commanded
neither Hebrew nor Greek. His genius was pastoral, moral, and administrative.
He translated the deep patristic inheritance---above all Augustine's vision of
interiority, charity, and divine sovereignty---into a clear, livable pattern
of Christian life that monks, bishops, missionaries, and rulers could actually
live by in a violent, impoverished world.

Looking straight at the primary evidence leaves no room for a plaster saint.
Gregory's letters reveal a man of fierce irritability under chronic gout,
occasionally prone to courtly flattery toward dangerous rulers such as
Brunehild and Phocas, willing to invoke imperial penalties against Donatists
and harsh pressure against stubborn pagan peasants, and convinced to his dying
day that the world was nearing its final hour. Yet beside those genuine limits
stands a rare, luminous integrity. He wept genuinely for the contemplative
stillness he had lost, yet never used longing for heaven as an excuse to
neglect a starving refugee, an exploited Sicilian tenant, a Jewish congregation
driven from its synagogue, or an unreached people in a distant northern island.
He defended the authority of Saint Peter's chair without tolerating the vanity
that diminishes brother bishops, insisting that his own honor lived only in
the honor of the universal Church.

Several historical questions remain open where the sixth-century record falls
silent. Neither Gregory's exact birth year nor the reputed link of his family
to the \latin{gens Anicia} can be proved from contemporary sources. The exact
date of his diaconal ordination and whether Benedict I or Pelagius II first
called him from Saint Andrew's remain plausible inference rather than certain
fact. How much Gregory knew of the murder of Maurice and his sons when he
drafted his congratulatory letters to Phocas in 603 cannot be settled with
certainty from the surviving \work{Registrum}. And while Gregory's hand in the
\latin{Hanc igitur}, the \latin{Pater noster}, and the Roman \latin{Kyrie} is
well attested, reconstructing the precise extent of his own liturgical texts
inside the later \work{Hadrianum} remains a specialist question.

What endures beyond every open question is the voice of the monk who stayed at
the helm of a battered ship until his strength ran out---teaching bishops that
the guidance of souls is the art of arts, praying in the middle of war that
God might dispose our days in his peace, and naming himself simply the servant
of the servants of God.
